\documentclass[a4paper,8pt]{extarticle}

% ---- Geometry: tight margins so everything fits on one A4 side ----
\usepackage[a4paper,top=8mm,bottom=8mm,left=9mm,right=9mm]{geometry}

\usepackage[T1]{fontenc}
\usepackage{lmodern}
\usepackage{array}
\usepackage{booktabs}
\usepackage{multicol}
\usepackage{xcolor}
\usepackage{enumitem}
\usepackage[hidelinks]{hyperref}
\usepackage{titlesec}

% ---- Kill vertical whitespace ----
\setlength{\parindent}{0pt}
\setlength{\parskip}{1.5pt}
\setlength{\columnsep}{6mm}
\setlength{\tabcolsep}{3pt}
\renewcommand{\arraystretch}{0.95}
\pagestyle{empty}

% ---- Compact section headings ----
\titleformat{\section}{\bfseries\small\color{black!85}}{\thesection.}{0.4em}{}
\titlespacing{\section}{0pt}{3pt}{1pt}

% ---- Helpers ----
\newcommand{\lic}[1]{\textbf{#1}}
\newcommand{\hd}[1]{\textbf{#1}}

\begin{document}
\footnotesize

\begin{center}
{\large\bfseries GymJam --- Copyright \& Legal Issues Report}\\[1pt]
{\scriptsize Every third-party resource (libraries, fonts, icons, map data, images) used in the GymJam digital touchpoint, its licence, and the
legal constraints a future team must respect before officially releasing the product. GymJam is an Expo / React Native app (iOS, Android, web) with a FastAPI + Supabase backend.}
\end{center}
\vspace{-2pt}
\hrule
\vspace{2pt}

\begin{multicols}{2}

\section{Third-party software libraries}

\textbf{Frontend --- JS/TS (\texttt{frontend/package.json}).} All \lic{MIT} unless noted:
react / react-dom / react-native (Meta); \texttt{expo} + every \texttt{expo-*} module (router, font, image, location, haptics, crypto, constants, linear-gradient, splash-screen, status-bar, symbols, system-ui, web-browser, apple-authentication, glass-effect, linking --- Expo/650 Industries);
\texttt{@react-navigation/*}; \texttt{@supabase/supabase-js}; \texttt{async-storage};
\texttt{react-native-maps} (\emph{map tiles behind it are not MIT --- see \S3});
Software Mansion libs (reanimated, gesture-handler, screens, svg, worklets);
safe-area-context; react-native-web; url-polyfill;
\texttt{@expo/vector-icons} (\emph{wrapper MIT; glyphs separately licensed --- \S2});
\texttt{@expo-google-fonts/plus-jakarta-sans} (\emph{font file is SIL OFL 1.1 --- \S2}).
Dev-only: jest / jest-expo, eslint / eslint-config-expo, \texttt{@types/*} (all MIT).
\texttt{typescript} (Microsoft) is \lic{Apache-2.0} --- preserve any \texttt{NOTICE}, state modifications, includes a patent grant (dev-only).

\smallskip
\textbf{Backend --- Python (\texttt{backend/requirements.txt}).}
fastapi (MIT); \texttt{uvicorn[standard]} (Encode, \lic{BSD-3-Clause} --- keep notice, don't use author's name to endorse); pydantic / pydantic-settings (MIT); supabase python client (MIT); httpx (Encode, \lic{BSD-3-Clause}); pytest (MIT, dev-only).
A local venv (\texttt{backend/venv/}) pulls in matplotlib + DejaVu/STIX fonts, but these are git-ignored, developer-only and \emph{not shipped} --- no release duty, but must not be committed.

\smallskip
\textbf{Position.} Every shipped dependency is \lic{permissive} (MIT, BSD-3-Clause, Apache-2.0). \textbf{No copyleft (GPL/AGPL/LGPL)}, so no duty to open-source GymJam's own code. The only practical duty is \textbf{attribution}: reproduce each library's copyright + licence text (auto-generated \texttt{licenses.txt} / in-app ``Open Source Licences'' screen).

\smallskip
\textbf{Hosting services (SaaS).} \emph{Supabase} (auth/DB/storage) acts as a \textbf{data processor} --- bound by its ToS + DPA; hosting region sets where personal data lives (\S\,GDPR). \emph{Expo/EAS} and \emph{Vercel} (web build) bound by their ToS. \emph{Overpass API / OpenStreetMap} (\texttt{fetch\_uk\_gyms.py}) --- data is \lic{ODbL}, with attribution + share-alike duties (\S3) and a fair-use policy on public mirrors.

\section{Fonts \& icons}

\begin{itemize}[leftmargin=1em,itemsep=0pt,topsep=1pt]
\item \textbf{Plus Jakarta Sans} (Tokotype) --- main app typeface + leaflet --- \lic{SIL OFL 1.1}: free to embed in commercial apps; \emph{cannot be sold standalone}; reserved font name must not be reused on a modified version.
\item \textbf{Inter} (Andersson) \& \textbf{Bricolage Grotesque} --- pitch-leaflet HTML via Google Fonts CDN --- \lic{SIL OFL 1.1} (as above).
\item \textbf{Material Icons / Symbols} (Google) --- in-app UI via \texttt{@expo/vector-icons} --- \lic{Apache-2.0}: free commercial use; preserve the notice.
\end{itemize}
The leaflet pulls fonts from the \textbf{Google Fonts CDN}; free for commercial use, but for a production web release consider \textbf{self-hosting} --- a German court has ruled dynamic CDN embedding can send user IPs to Google and breach GDPR.

\section{Map \& geographic data (key external resource)}

\textbf{react-native-maps} uses \texttt{PROVIDER\_DEFAULT} (\texttt{FullMap.tsx}, \texttt{ProfileMap.tsx}) $\Rightarrow$ \textbf{Apple Maps on iOS}, \textbf{Google Maps on Android}.
\emph{Apple MapKit} --- governed by the Apple Developer Program Licence; no separate fee.
\emph{Google Maps Platform} --- billable above the free tier; needs an API key + billing for a public release; data/imagery may not be scraped or used outside a Google map.

\smallskip
\textbf{Gym dataset} (\texttt{gyms\_seed.sql} via \texttt{fetch\_uk\_gyms.py}) is derived from \textbf{OpenStreetMap / Overpass} under the \lic{Open Database Licence (ODbL) v1.0}. Duties: (a) \textbf{attribution} --- show ``© OpenStreetMap contributors'' wherever the data appears; (b) \textbf{share-alike} --- any publicly distributed derived database must itself be ODbL; (c) keep the data open. \emph{For release: add a visible OSM credit on/near the map screen.}

\smallskip
\textbf{Images \& original assets.} App icons/splash/adaptive icons, mascot pixel-art SVGs (\texttt{fit, scrawny, mogger, buff, otter}), leaflet imagery and the ``GJ'' logotype are all \textbf{group-owned original work} (brand subject to trademark clearance). No third-party photos / stock / clip-art are embedded --- the image-rights position is clean; track licence + source for any future stock additions.

\section{Wider legal implications}

\textbf{Data protection (UK GDPR / DPA 2018 / EU GDPR).} GymJam collects \textbf{personal data}: email + password (Supabase auth), precise \textbf{device location} (\texttt{expo-location}, maps), social graph, and behavioural data (check-ins, streaks, pledges). \texttt{Login.tsx} already shows ``you agree to GymJam's Terms of Service and Privacy Policy'' --- \textbf{neither document exists and both must be written}. Obligations: a lawful basis for processing, \textbf{explicit consent / justification for location}, data-subject rights (access, deletion, portability), data minimisation, breach notification, and a \textbf{DPA with Supabase}. As Supabase may host in the \textbf{US}, an international-transfer mechanism (SCCs) is needed; password handling mandates appropriate security.

\smallskip
\textbf{Financial / gambling regulation (highest risk).} The ``pot''/pledge mechanic (\texttt{services/pot.py}, pledge flows) lets users stake value on showing up. If \textbf{real money} is ever involved this may engage \textbf{gambling law (UK Gambling Act 2005)} and/or \textbf{payment-services / FCA} rules (client funds, e-money). Get legal advice before launch; keeping pledges to non-monetary points/ELO substantially reduces risk.

\smallskip
\textbf{Health \& liability.} As a fitness product encouraging gym attendance, GymJam should carry a \textbf{medical/injury disclaimer} (``consult a doctor before exercising'') and limit liability in its ToS.

\smallskip
\textbf{App-store policies.} A release must satisfy \textbf{Apple App Store Review Guidelines} and \textbf{Google Play policies}: published privacy policy, accurate permission strings (location etc.), account-deletion support, and --- since \texttt{expo-apple-authentication} is a dependency --- Apple's ``Sign in with Apple'' parity rule if any third-party social login is offered.

\smallskip
\textbf{Trademark \& brand.} Clear the name \textbf{``GymJam''} against the UK trademark register before launch. The dataset references gym-chain names (PureGym, The Gym Group, Virgin Active); using them \textbf{as factual location names} is generally permissible \emph{nominative use}, but the app must \textbf{not} use their logos or imply partnership/endorsement.

\smallskip
\textbf{Open-source attribution.} Though all licences are permissive, MIT/BSD/Apache all require their notices to be reproduced. Ship an \textbf{``Open Source Licences'' acknowledgements screen} (or bundled \texttt{NOTICE}/\texttt{licenses.txt}) listing every \S1 dependency, the \S2 fonts, and the \textbf{``© OpenStreetMap contributors''} credit from \S3.

\end{multicols}
\end{document}
